% Noether P06 Korean producer draft — UNCHECKED
% T16-U224; ED0002 lines5810--5812; source 637 LF B / DA5E387EF877B778543D8FC1EDF748629A80C509D9A729DFF450F6AC64606FCE
% Pointer v009 / ED0002: B06BE3530D9CF2E82B56FDBA7FE41D5D044DF2425DFA2A059D4939EAA2F7A6C2 / C9A125167ACB33D914EE4374B65AE7CDF0052F568371B8B77B720EA178ABF0E3
% Translation only; no review/build/render/approval. Stronger lemma derivation, A-coefficient function class, sufficient condition, and unit resultant remain held.

정수계수 정기저의 존재증명을 옮기기 위해, \S\ 11의 보조정리가 다음과 같은 더 강한 형태를 갖는다는 점에 유의한다. 이는 \S\ 11의 방정식 (5), (6)과 $A_1(\lambda)\rcdots A_\tau(\lambda)$가 $\lambda_1\rcdots\lambda_\rho$의 정수계수 다항함수라는 사실에서 따른다.

\emph{$x_1\rcdots x_\rho$의 임의의 모든 정수계수 다항식이 $\rho$개의 정수계수 다항식에 대하여 대수적으로 정수적이고 정수계수적으로 종속되기 위한 충분조건은, 이 다항식들의 최고차항들의 결과식이 $[\Omega]$의 한 단원인 것이다 --- $R_0=\eps$.}
